%% tcbpagefootnote.tex -- documentation for the tcbpagefootnote package
%%
%% Copyright (c) 2026 Munehiro Yamamoto <munepixyz@gmail.com>
%% License: LPPL 1.3c or later
%%
%% Build: pdflatex tcbpagefootnote.tex (twice)
\documentclass{article}
\usepackage[margin=28mm]{geometry}
\usepackage{xcolor}
\usepackage{tcbpagefootnote}
\usepackage[colorlinks, allcolors=blue!50!black]{hyperref}
\newcommand*\tpfversion{v1.0 (2026/10/25)}
\newcommand*\cs[1]{\texttt{\char`\\#1}}
\newcommand*\pkg[1]{\textsf{#1}}
\newcommand*\key[1]{\texttt{#1}}
\newcommand*\meta[1]{$\langle$\textit{#1}$\rangle$}
\tcbset{colframe=black!60, colback=black!3, fonttitle=\bfseries}

\title{The \pkg{tcbpagefootnote} package\\
  \large Page footnotes from inside tcolorbox}
\author{Munehiro Yamamoto\\
  \url{https://github.com/munepi/tcbpagefootnote}}
\date{\tpfversion}

\begin{document}
\maketitle

\section{Introduction}

\pkg{tcolorbox} typesets the content of a box like a \texttt{minipage}:
a \cs{footnote} inside the box is numbered with the \texttt{mpfootnote}
counter (a, b, \dots) and its text is put at the bottom of the box.
This is often what you want for a box that stands apart from the text,
but not for a box that is part of the running text, such as a note or
a definition in a book: there, the footnotes should be numbered along
with those of the body text and put at the bottom of the page.

The \pkg{tcbpagefootnote} package adds the \pkg{tcolorbox} key
\key{page footnote} for this.  The two boxes below differ only in this
key.

\begin{tcolorbox}[title=An ordinary box]
A footnote in an ordinary box\footnote{This footnote is put at the
  bottom of the box.} is a footnote of the box.
\end{tcolorbox}

\begin{tcolorbox}[title=A box with \key{page footnote}, page footnote]
A footnote in a box with \key{page footnote}\footnote{This footnote is
  put at the bottom of the page.} is a footnote of the page, numbered
along with the footnotes of the body text.
\end{tcolorbox}

\paragraph{Supported formats/engines}
pdf\LaTeX, Xe\LaTeX, Lua\LaTeX{} (with or without Lua\TeX-ja),
p\LaTeX{} and up\LaTeX{} (with dvipdfmx).  The package requires
\pkg{tcolorbox} and loads its \texttt{hooks} library.

\section{Usage}

\begin{verbatim}
\usepackage{tcbpagefootnote}

\begin{tcolorbox}[title=Note, page footnote]
  Text\footnote{A page footnote.}.
\end{tcolorbox}

\newtcolorbox{note}[1][]{colback=blue!5, page footnote, #1}
\tcbset{page footnote}  % for every box
\end{verbatim}

\noindent
The key acts on the footnotes in the title, the upper part and the lower
part of the box:
\begin{itemize}
\item the marks are made with the \texttt{footnote} counter and
  \cs{thefootnote}, as in the body text;
\item the texts are collected while the box is typeset and are put into
  the ordinary footnote insertion right after the box, so that they go
  to the bottom of the page where the box is.  If the box does not fit
  on the current page and moves to the next one, its footnotes move
  with it.
\end{itemize}
\cs{footnote}\texttt{[\meta{number}]}, \cs{footnotemark} with
\cs{footnotetext}, \cs{tcbox}, and boxes with \key{page footnote} inside
each other work as usual.  An ordinary box inside a box with
\key{page footnote} keeps its own footnotes.  With \pkg{hyperref}, each
footnote mark links to its own footnote text.

\section{Errors and warnings}

\begin{description}
\item[\texttt{The key 'page footnote' cannot be used with 'float'.}]
  A floating box (\key{float}) is not followed by the code that puts the
  footnotes.
\item[\texttt{The 'after' code of this box no longer flushes the page
  footnotes.}] The footnotes are put by the \key{after pre} code of the
  \texttt{hooks} library.  Keys that set the whole \texttt{after} code
  (\key{after}, \key{after skip}, \key{nobeforeafter}, \dots) remove it
  when they come after \key{page footnote}: give them before it, e.g.\
  \verb|[after skip=1pt, page footnote]|.
\item[\texttt{The tcolorbox key 'page footnote' is already defined.}]
  Another package defines a key of the same name; \pkg{tcbpagefootnote}
  does not redefine it.
\item[\texttt{This box is broken across pages.} (warning)] See below.
\end{description}

\section{Limitations}

\begin{itemize}
\item A \key{breakable} box that is broken across pages puts the
  footnote texts of all its parts on the page where the box ends: the
  content of a breakable box is typeset as a whole before it is split,
  and footnote insertions do not travel with the split parts.  The
  package issues a warning for such a box when it has footnotes.  The
  footnotes of a breakable box that is not broken are put as usual.
\item \key{page footnote} cannot be used with \key{float}.
\item Vertical writing (tate) is outside the scope of this version.  With
  (u)p\LaTeX{} and dvipdfmx, \pkg{tcolorbox} itself does not work in
  tate, since the dvipdfmx driver of TikZ/PGF does not support it;
  this is to be handled by a separate package.
\end{itemize}

\section{Tests}

\verb|make test| compiles the test files with pdf\LaTeX, Xe\LaTeX,
Lua\LaTeX, p\LaTeX{} and up\LaTeX{} (+ dvipdfmx) and checks the log.
For every footnote, \texttt{test-common.tex} records in the
\texttt{.aux} file at shipout the page and the mark of both the footnote
mark and the footnote text, and whether the text is set in the full
\cs{columnwidth} of the page (a page footnote) or in a box; when the
\texttt{.aux} file is read back at \verb|\end{document}|, it checks
that every text is typeset exactly once, with its own mark, on the
expected page (\texttt{TPF-RESULT OK}).

\begin{center}
\begin{tabular}{ll}
  \hline
  Test & What it covers\\
  \hline
  \texttt{test-basic.tex}     & title/upper/lower, nesting, \cs{tcbox},
                                a box moved to the next page\\
  \texttt{test-hyperref.tex}  & \texttt{test-basic.tex} with
                                \pkg{hyperref}: no duplicate anchors\\
  \texttt{test-breakable.tex} & broken and unbroken breakable boxes,
                                the warning\\
  \texttt{test-japanese.tex}  & Japanese classes (ltjbook, jbook,
                                ujbook)\\
  \texttt{test-errors.tex}    & the error messages\\
  \hline
\end{tabular}
\end{center}

\section{History}

\begin{description}
\item[v1.0 (2026/10/25)] First public release.
\end{description}

\section{License}

This package is licensed under the conditions of the \LaTeX{} Project
Public License, either version~1.3c of this license or any later
version.

\medskip
\noindent
Copyright \copyright{} 2026 Munehiro Yamamoto

\end{document}
